\TUDoption{abstract}{section, multiple}

% TODO: Replace the placeholders with your actual abstracts

\newcommand*{\Abstreng}{Estimating high-resolution human textures from a single image remains a core challenge in 3D human modeling, this problem is crucial for creating photorealistic 3D virtual humans in applications ranging from cinema industry, game industry to Virtual Reality(VR),Augmented Reality(AR) development. While prior works such as SMPLitex\cite{casas2023smplitex} have explored the use of diffusion models\cite{rombach2021highresolution} for UV texture\cite{10.1145/566654.566590} completion, their output resolution (512²) and data scale remain limited, restricting downstream application and generalization. In this thesis, I propose High-Resolution Human Texture estimation(HRHumanTex), a new framework that enables 1024²-resolution UV texture estimation from a single image by integrating state-of-the-art super-resolution (SR) and Latent Diffusion Model inpainting strategies into a unified pipeline.

First, I construct enhanced UV datasets by applying seven state-of-the-art(SOTA) super resolution(SR) models—including RealESRGAN\cite{wang2021realesrgan}, SRFormer\cite{zhou2023srformer}, DiffBIR\cite{lin2024diffbir}, SwinIR\cite{liang2021swinir}, BSRGAN\cite{zhang2021designing}, StableSR\cite{wang2024exploiting} and RCAN\cite{zhang2018rcan}—on SMPLitex textures, and benchmark their effectiveness on human UV texturemap enhancement. Next, I fine-tune Stable Diffusion v1.5 with DreamBooth\cite{ruiz2023dreamboothfinetuningtexttoimage} on each SR-augmented dataset, resulting in multiple specialized inpainting models. Given a single RGB image, HRHumanTex computes a masked partial texture using DensePose\cite{wu2019detectron2} and SGHM\cite{chen2022sghm}, and completes the full texture with one of the trained diffusion inpainting models.

To evaluate both SR and inpainting stages, I introduce a multi-level evaluation suite: traditional metrics including Peak Signal-to-Noise Ratio (PSNR), Structural Similarity Index (SSIM), Learned Perceptual Image Patch Similarity (LPIPS), Fréchet Inception Distance (FID), rendering based SSIM and LPIPS evaluation by estimating pose and camera parameters with ROMP\cite{ROMP}, and rendering SMPL model with uv texture then evaluate them with ground truth images. Differentiable rendering\cite{ravi2020pytorch3d} based uv texture optimization also help to get better texture quality. Extensive experiments on DeepFashion\cite{liuLQWTcvpr16DeepFashion} and SHHQ\cite{fu2022styleganhuman} show that HRHumanTex generates sharper, more realistic and more consistent UV textures compared to SMPLitex baseline methods. 

Overall, this work demonstrates how scaling up both data quality and model resolution leads to significant gains in single-image 3D texture estimation, providing new insights for texture synthesis, virtual try-on(VTO), and human digitization.}

\newcommand*{\Abstrger}{Die Schätzung hochauflösender menschlicher Texturen aus einem einzigen Bild bleibt eine zentrale Herausforderung in der 3D-Humanmodellierung. Dieses Problem ist entscheidend für die Erstellung fotorealistischer virtueller Menschen in Anwendungen wie Filmindustrie, Spieleentwicklung sowie Virtual Reality (VR) und Augmented Reality (AR). Während frühere Arbeiten wie SMPLitex~\cite{casas2023smplitex} den Einsatz von Diffusionsmodellen~\cite{rombach2021highresolution} zur Vervollständigung von UV-Texturen~\cite{10.1145/566654.566590} untersucht haben, sind deren Ausgaberesolution (512²) und Datengröße begrenzt, was die Anwendbarkeit und Generalisierbarkeit einschränkt.

In dieser Arbeit schlage ich \textbf{HRHumanTex} (High-Resolution Human Texture Estimation) vor – ein neues Framework zur Erzeugung von UV-Texturen in $1024^2$-Auflösung aus einem einzelnen Bild. Dies wird durch die Integration von modernen Super-Resolution (SR)-Methoden und Latent-Diffusionsmodellen zur Inpainting in eine einheitliche Pipeline erreicht.

Zunächst konstruiere ich erweiterte UV-Datensätze, indem ich sieben aktuelle SR-Modelle – darunter RealESRGAN~\cite{wang2021realesrgan}, SRFormer~\cite{zhou2023srformer}, DiffBIR~\cite{lin2024diffbir}, SwinIR~\cite{liang2021swinir}, BSRGAN~\cite{zhang2021designing}, StableSR~\cite{wang2024exploiting} und RCAN~\cite{zhang2018rcan} – auf SMPLitex-Texturen anwende und deren Effektivität auf UV-Texturen von Menschen vergleiche. Anschließend feintune ich Stable Diffusion v1.5 mit DreamBooth~\cite{ruiz2023dreamboothfinetuningtexttoimage} auf jedem SR-augmentierten Datensatz und erhalte spezialisierte Inpainting-Modelle. Ausgehend von einem RGB-Bild erzeugt HRHumanTex eine partielle UV-Texturmaske mit DensePose~\cite{wu2019detectron2} und SGHM~\cite{chen2022sghm}, welche anschließend durch ein trainiertes Diffusionsmodell vervollständigt wird.

Zur Bewertung der SR- und Inpainting-Stufen führe ich ein mehrstufiges Evaluierungssystem ein: klassische Metriken wie Peak Signal-to-Noise Ratio (PSNR), Structural Similarity Index (SSIM), Learned Perceptual Image Patch Similarity (LPIPS) und Fréchet Inception Distance (FID), sowie eine renderbasierte SSIM/LPIPS-Bewertung. Letztere basiert auf Posen- und Kameraschätzung mittels ROMP~\cite{ROMP} und einer differentiell renderbaren Projektion der SMPL-Texturen~\cite{ravi2020pytorch3d}. Darüber hinaus nutze ich differentielles Rendering zur weiteren Optimierung der UV-Texturqualität.

Umfangreiche Experimente auf DeepFashion~\cite{liuLQWTcvpr16DeepFashion} und SHHQ~\cite{fu2022styleganhuman} zeigen, dass HRHumanTex im Vergleich zu SMPLitex deutlich schärfere, realistischere und konsistentere UV-Texturen erzeugt.

Insgesamt zeigt diese Arbeit, dass eine Skalierung sowohl der Datenqualität als auch der Modellauflösung signifikante Fortschritte bei der Texturschätzung aus einem einzigen Bild ermöglicht – mit Potenzial für Anwendungen wie Texture Synthesis, Virtual Try-On (VTO) und Menschendigitalisierung.}

\begin{abstract}[pagestyle=empty.tudheadings]
	\iflanguage{ngerman}
		{\Abstrger}
		{\Abstreng}
	
	\iflanguage{ngerman}
	{\nextabstract[english]}
	{\nextabstract[ngerman]}
	
	\iflanguage{ngerman}{
		\Abstrger}{
		\Abstreng
	}

\end{abstract}
